\documentclass[12pt,a4paper,openany]{book}
\usepackage{verbatim}
\usepackage[UKenglish]{babel}
\usepackage[UKenglish]{isodate}
\usepackage[T1]{fontenc}
\usepackage{lmodern}
\usepackage[margin=1in]{geometry}
\cleanlookdateon
\usepackage{epsfig}
\usepackage{graphicx}
\usepackage[margin=15pt,font={footnotesize},labelfont=bf,format=hang]{caption}
\usepackage[margin=15pt,font={footnotesize},labelfont=bf,format=hang]{subcaption}
\usepackage{bm}
\usepackage{amsbsy}
\usepackage{amssymb}
\usepackage{amsmath}
\usepackage{fontspec}
\setsansfont{Fira Sans}

\usepackage{isomath}

\usepackage{mathtools}
\usepackage{enumitem}
\usepackage{natbib}
\usepackage[pdfencoding=auto,pdftitle={MATH3171 Mathematical Biology III Course Notes},pdflang={en-GB}]{hyperref}
\usepackage[dvipsnames]{xcolor}
\usepackage[ruled,vlined,linesnumbered]{algorithm2e}
\usepackage{cleveref}
\usepackage{nameref}
\usepackage{xpatch}
\xpretocmd{\algorithm}{\hsize=\linewidth}{}{}
\usepackage[normalem]{ulem}
\usepackage{booktabs}
\usepackage{soul}
\usepackage{siunitx}
\usepackage{color}
\usepackage{tikz}
\usetikzlibrary{shapes,arrows,calc}
\usetikzlibrary{positioning}
\usepackage{authblk}
\usepackage{etoolbox}
\usepackage[most]{tcolorbox}
\tcbuselibrary{breakable}
\usetikzlibrary{arrows.meta}
\usepackage{circledsteps}
\usepackage{microtype}


\usepackage{pgfplots}

\usepackage{subfiles}

\usepackage{newunicodechar}
\UndeclareTextCommand{\textellipsis}{TU}
\newunicodechar{…}{\ldots}

\usepackage[raggedright]{titlesec}

\titleformat{\chapter}[display]
  {\normalsize \huge \color{black}}%
  {\LARGE \fontspec{Fira Sans Medium} \addfontfeature{LetterSpace=15} \MakeUppercase{\chaptertitlename{}\, \thechapter}
  }%
  {0ex}%
  {\fontspec[Ligatures=TeX]{Fira Sans Medium}\huge\color{adam-purple}\filright}%

\titleformat{\part}[display]
  {\normalsize \huge \color{black}}%
  {\centering \LARGE \fontspec{Fira Sans Medium} \addfontfeature{LetterSpace=15} \MakeUppercase{\partname{}\, \thepart}
	}%
  {0ex}%
  {\centering \fontspec[Ligatures=TeX]{Fira Sans Medium}\huge\color{adam-purple}}%

\usepackage{fancyhdr}
	\pagestyle{fancy}
	\renewcommand{\headrulewidth}{0pt}
	\cfoot{}
	\renewcommand\sectionmark[1]{\markright{\thesection\; #1}}
	\fancyhead[LO]{\small\it\lowercase\nouppercase\selectfont\rightmark}
	\fancyhead[RO]{{\thepage}}
	\fancyhead[RE]{\small\it\lowercase\nouppercase\selectfont\rightmark}
	\fancyhead[LE]{{\thepage}}
	\setlength{\headheight}{15pt}
	\fancypagestyle{blank}{\fancyhf{}\fancyfoot[C]{}\renewcommand{\headrulewidth}{0pt}}

% Vectors, quaternions etc.
\renewcommand{\v}[1]{\mathbfit{#1}}      % Generic vector
\newcommand{\vc}[1]{\bm{\mathcal{#1}}}      % vector mathcal
\newcommand{\uv}[1]{\widehat{\mathbfit{#1}}}      % Generic unit vector
\newcommand{\q}[1]{\mathsfbfit{#1}}        % Generic quaternion
\renewcommand{\t}[1]{\mathsfbfit{#1}}	 % Generic tensor
\newcommand{\m}[1]{\mathsfbfit{#1}}	 % Generic matrix
\newcommand{\vu}[1]{\mathbf{#1}}         % Upright vector (for zero vector)
\newcommand{\qu}[1]{\mathbf{#1}}       % Upright quaternion (for zero quaternion)
\newcommand{\tu}[1]{\bm{\mathsf{#1}}}	 % Upright tensor (for zero tensor)

%Vector operators
\renewcommand{\d}{\mathrm{d}}
\def \p {\partial}
\newcommand{\del}{\nabla}
\newcommand{\grad}{\nabla}
\newcommand{\vdel}{\boldsymbol{\nabla}}
\newcommand{\vgrad}{\boldsymbol{\nabla}}
\newcommand{\vnabla}{\bm{\nabla}}
\newcommand{\cross}{{\times}}
\newcommand{\Lap}{{\nabla^2}}

%Hats, tildes and bars that fit
\renewcommand{\hat}{\widehat}
\renewcommand{\tilde}{\widetilde}
\renewcommand{\bar}{\overline}

% Differentiating
\newcommand{\fd}[2]{\mathchoice{\frac{\d #1}{\d #2}}{\d #1/\d #2}{\d #1/\d #2}{\d #1/\d #2}}
\newcommand{\fdd}[2]{\mathchoice{\frac{\d^2 #1}{\d #2^2}}{\d^2 #1/\d #2^2}{\d^2 #1/\d #2^2}{\d^2 #1/\d #2^2}}
\newcommand{\fddd}[2]{\mathchoice{\frac{\d^3 #1}{\d #2^3}}{\d^3 #1/\d #2^3}{\d^3 #1/\d #2^3}{\d^3 #1/\d #2^3}}
\newcommand{\fdddd}[2]{\mathchoice{\frac{\d^4 #1}{\d #2^4}}{\d^4 #1/\d #2^4}{\d^4 #1/\d #2^4}{\d^4 #1/\d #2^4}}
\newcommand{\fdn}[3]{\mathchoice{\frac{\d^{#1} #2}{\d #3^{#1}}}{\d^{#1} #2/\d #3^{#1}}{\d^{#1} #2/\d #3^{#1}}{\d^{#1} #2/\d #3^{#1}}}
\newcommand{\pd}[2]{\mathchoice{\frac{\p #1}{\p #2}}{\p #1/\p #2}{\p #1/\p #2}{\p #1/\p #2}}
\newcommand{\pdd}[2]{\mathchoice{\frac{\p^2 #1}{\p #2^2}}{\p^2 #1/\p #2^2}{\p^2 #1/\p #2^2}{\p^2 #1/\p #2^2}}
\newcommand{\pddmixed}[3]{\mathchoice{\frac{\p^2 #1}{\p #2 \p #3}}{\p^2 #1 /\p #2 \p #3}{\p^2 #1 /\p #2 \p #3}{\p^2 #1 /\p #2 \p #3}}
\newcommand{\pddd}[2]{\frac{\p^3 #1}{\p #2^3}}
\newcommand{\pdn}[3]{\mathchoice{\frac{\p^{#1} #2}{\p #3^{#1}}}{\p^{#1} #2/\p #3^{#1}}{\p^{#1} #2/\p #3^{#1}}{\p^{#1} #2/\p #3^{#1}}}

% Misc
\DeclareMathSymbol{\Delta}{\mathalpha}{operators}{1} % Keeps Delta upright
\DeclareMathSymbol{\Gamma}{\mathalpha}{operators}{0} % Keeps Gamma upright
\newcommand{\cond}{\kappa}
\newcommand{\diag}{\operatorname{diag}}
\newcommand{\tr}{\operatorname{tr}}
\newcommand{\erf}{\operatorname{erf}}
\newcommand{\order}{\mathcal{O}}
\newcommand{\Reals}{\mathbb{R}}
\newcommand{\twopartdef}[4]{
	\left\{
		\begin{array}{ll}
			#1 & \mbox{} #2 \\[6pt]
			#3 & \mbox{} #4
		\end{array}
	\right.
}

% Misc
\newcommand{\e}{\mathrm{e}}
\renewcommand{\i}{\mathrm{i}}

\newcommand{\dx}{\,\mathrm{d}x}
\newcommand{\dy}{\,\mathrm{d}y}
\renewcommand{\epsilon}{\varepsilon}
\newcommand{\ep}{\epsilon}
\renewcommand{\Re}{\operatorname{Re}}
\newcommand{\Four}{\mathcal{F}}
\newcommand{\Lin}{\mathcal{L}}


\newcommand{\kmi}{\widehat{\v{k}\cdot\v{x}_i}}
\newcommand{\norm}[1]{\left\vert #1 \right\vert}
\newcommand{\lb}{\left \langle}
\newcommand{\rb}{\right \rangle}
%\newcommand{\ks}{|\v{k}|}
\newcommand{\ks}{\rho_k}


	%Column vectors
%	\makeatletter
%\newcommand\rcvector[2][\\]{\ensuremath{%
%  \global\def\rc@delim{#1}%
%    \negthinspace\begin{pmatrix}
%      \rc@vector #2,\relax\noexpand\@eolst%
%    \end{pmatrix}}}
%\def\rc@vector #1,#2\@eolst{%
%  \ifx\relax#2\relax
%    #1
%  \else
%    #1\rc@delim
%    \rc@vector #2\@eolst%
%  \fi}
%\makeatother
\newcommand{\rcvector}[1]{
\begin{pmatrix}
#1
\end{pmatrix}
}

\newcommand{\colvec}{\rcvector}
\newcommand{\rowvec}{\rcvector}
\renewcommand{\binom}{\rcvector}

\setlength{\parindent}{0ex}
\setlength{\parskip}{2ex}

\newcommand{\deriv}[2]{\frac{\d#1}{\d #2}}
\newcommand{\derivn}[3]{\frac{\d^{#3} #1}{\d #2^{#3}}}
\newcommand{\pder}[2]{\frac{\partial #1}{\partial #2}}
\newcommand{\pdern}[3]{\frac{\partial^{#3} #1}{\partial #2^{#3}}}
\newcommand{\mvec}[2]{\left(\begin{array}{c}#1\\#2\end{array}\right)}
\newcommand{\mat}[4]{\begin{pmatrix}#1 & #2 \\ #3 & #4 \end{pmatrix}}
\newcommand{\im}{\i }
%\newcommand{\matrix}[1]{\begin{pmatrix}#1

\newcommand{\adam}[1]{{\color{red}#1}}

\newtheorem{definition}{Definition}[chapter]

\definecolor{adam-purple}{HTML}{5f466d}
\definecolor{question-purple}{HTML}{f6f1fa}
\definecolor{nem-red}{HTML}{FFEDF0}
\definecolor{nem-dark-red}{HTML}{d21f3c}
\definecolor{link-blue}{HTML}{0e4d92}
\definecolor{info-dark-brown}{HTML}{7a4501}
\definecolor{info-brown}{HTML}{FFF4E5}
\definecolor{reading-dark-blue}{HTML}{0e4d92}
\definecolor{reading-blue}{HTML}{e8f1ff}
\definecolor{C0}{HTML}{1f77b4}
\definecolor{C1}{HTML}{ff7f0e}
\definecolor{C2}{HTML}{2ca02c}
\definecolor{C3}{HTML}{d62728}
\definecolor{C4}{HTML}{9467bd}
\definecolor{C5}{HTML}{8c564b}
\definecolor{C6}{HTML}{e377c2}
\newcommand{\colone}[1]{{\color{C0}#1}}
\newcommand{\coltwo}[1]{{\color{C1}#1}}
% \newenvironment{do-it-yourself}{\begin{tcolorbox}[colback=question-purple,colframe=adam-purple,title={\sf\small{Have a go!}\vspace{-0.5mm}}]}{\end{tcolorbox}}
% \newenvironment{info}{\begin{tcolorbox}[colback=info-brown,colframe=info-dark-brown,title={\sf\small{Information}}]}{\end{tcolorbox}}
% \newenvironment{nem}{\begin{tcolorbox}[colback=nem-red,colframe=nem-dark-red,breakable,enhanced,parbox=false,title={\sf\small{Non-examinable}}]}{\end{tcolorbox}}
% \newenvironment{extra-reading}{\begin{tcolorbox}[colback=reading-blue,colframe=reading-dark-blue,title={\sf\small{Extra reading}\vspace{-0.5mm}}]}{\end{tcolorbox}}

\newenvironment{colour-box}[3]{\begin{tcolorbox}[
	enhanced,
	colback=#2,
	colframe=#1,
	colbacktitle=#2,
	%frame hidden,
	boxrule=0pt,
	leftrule=2pt,
	toptitle=2.5mm,
	bottomtitle=-.5mm,
	bottom=3mm,
  sharp corners=west,
  frame code={
    \fill[#1]
      ([xshift=2pt]frame.north west) --
      ([xshift=2pt]frame.south west) arc
      [radius=2pt,start angle=-90,end angle=-180] --
      ([yshift=-2pt]frame.north west) arc
      [radius=2pt,start angle=180,end angle=90]--
      cycle;
  },
  	fonttitle={\sf\small\setmainfont{Fira Sans Medium}},
	title={\color{#1}#3}
]}{\end{tcolorbox}}

\newenvironment{colour-box-breakable}[3]{
\setlist[itemize]{topsep=-\parskip,before=\vspace{-2mm},itemsep=3pt}
%\setlength{\parskip}{0pt}
\begin{tcolorbox}[
	enhanced,
    breakable,
	colback=#2,
	colframe=#1,
	colbacktitle=#2,
	frame hidden,
    boxrule=0pt,
    leftrule=2pt,
	toptitle=2.5mm,
	bottomtitle=-.5mm,
    pad before break=1pt,
	pad after break=6pt,
	parbox=false,
	bottom=10pt,
	%grow to right by=0.5mm,
    %grow to left by=-0.2mm,
	sharp corners=west,
  	fonttitle={\sf\small\setmainfont{Fira Sans Medium}},
	title={\color{#1}#3},
  	overlay first={
	  \fill[#1]
      ([xshift=2pt]frame.north west) --
      ([xshift=2pt]frame.south west) --
      ([xshift=0pt]frame.south west) --
      ([yshift=-2pt]frame.north west) arc
      [radius=2pt,start angle=180,end angle=90]--
      cycle;
    },
     overlay unbroken={
	  \fill[#1]
      ([xshift=2pt]frame.north west) --
      ([xshift=2pt]frame.south west) arc
      [radius=2pt,start angle=-90,end angle=-180] --
      ([yshift=-2pt]frame.north west) arc
      [radius=2pt,start angle=180,end angle=90]--
      cycle;
    },
  	overlay last={
	  \fill[#1]
      ([xshift=2pt]frame.north west) --
      ([xshift=2pt]frame.south west) arc
      [radius=2pt,start angle=-90,end angle=-180] --
      ([yshift=-2pt]frame.north west) --
      ([yshift=0pt]frame.north west) --
      cycle;
    },
]}{\end{tcolorbox}}

\newenvironment{do-it-yourself}{
	\begin{colour-box}
		{adam-purple}
		{question-purple}
		{\raisebox{0.5pt}{\scalebox{1.2}{$\blacktriangleright$}}\;Have a go!}
}{\end{colour-box}}

\newenvironment{info}{
	\begin{colour-box}
		{info-dark-brown}
		{info-brown}
		{\raisebox{0pt}{\scalebox{1}{$\blacksquare$}}\,\,Information}
}{\end{colour-box}}

\newenvironment{nem}{
	\begin{colour-box-breakable}
		{nem-dark-red}
		{nem-red}
		{\raisebox{0pt}{\scalebox{1.1}{$\blacklozenge$}}\;Non-examinable}
}{\end{colour-box-breakable}}

\newenvironment{extra-reading}{
	\begin{colour-box-breakable}
		{reading-dark-blue}
		{reading-blue}
		{\raisebox{1pt}{\scalebox{1}{$\bigstar$}}\;Extra reading}
}{\end{colour-box-breakable}}

\hypersetup{
    colorlinks,
    linkcolor={link-blue},
    citecolor={link-blue},
    urlcolor={link-blue}
}
\pgfkeys{/csteps/inner color=white}
\pgfkeys{/csteps/outer color=adam-purple}
\pgfkeys{/csteps/fill color=adam-purple}
\pgfkeys{/csteps/inner ysep=6pt}
\pgfkeys{/csteps/inner xsep=6pt}
\newcommand{\StepOneIcon}{\Circled{\smaller\sf 1}}
\newcommand{\StepTwoIcon}{\Circled{\smaller\sf\hspace{0.1ex}2}}
\newcommand{\StepThreeIcon}{\Circled{\smaller\sf\hspace{0.1ex}3}}
\newcommand{\StepFourIcon}{\Circled{\smaller\sf 4}}
\newcommand{\StepOne}{\StepOneIcon\;\;}
\newcommand{\StepTwo}{\StepTwoIcon\;\;}
\newcommand{\StepThree}{\StepThreeIcon\;\;}
\newcommand{\StepFour}{\StepFourIcon\;\;}

\newcommand{\flowchart}[1]{{\sf\ref{#1}.} \\ \nameref*{#1}}

\providecommand{\mathdefault}[1]{#1}